%======================= DESENVOLVIMENTO =======================
\section{Desenvolvimento}
\label{sec:desenvolvimento}

% Escrita a partir do código. Todos os números desta seção são reais
% e foram medidos. Faltam apenas as capturas de tela, marcadas com
% [FIGURA].

Esta seção descreve o artefato construído e, sobretudo, as decisões de
projeto que o determinaram. A ênfase recai sobre as decisões porque é
nelas que reside a contribuição de engenharia do trabalho: a maior parte
delas não decorre de preferência, mas de restrições concretas ---
perceptivas, de desempenho ou de plataforma --- que só se manifestaram
durante a construção.

\subsection{Tecnologias e organização}

O protótipo é uma aplicação web estática, executável diretamente no
navegador, sem instalação. A camada tridimensional utiliza o A-Frame
(versão 1.7.0), arcabouço declarativo construído sobre o Three.js que
implementa a especificação WebXR. Isso permite que a mesma página
funcione em três modos de uso --- computador, celular e óculos de
Realidade Virtual --- sem base de código separada para cada um.

A opção por HTML, CSS e JavaScript sem arcabouços de interface foi
deliberada, por três razões: reduz a dependência de terceiros, elimina
etapa de compilação (a aplicação é publicada copiando arquivos), e
mantém o código legível como artefato acadêmico --- um leitor que
conheça JavaScript consegue acompanhá-lo sem conhecer um arcabouço
específico.

\begin{table}[htbp]
\centering
\caption{Organização do código-fonte da experiência.}
\label{tab:arquivos}
\begin{tabular}{llr}
\toprule
\textbf{Arquivo} & \textbf{Responsabilidade} & \textbf{Linhas} \\
\midrule
\texttt{dados.js}        & os cinco animais e os cinco níveis            & 321 \\
\texttt{cenarios.js}     & monta os quatro ambientes e os esconderijos   & 597 \\
\texttt{experiencia.js}  & lê a escolha, monta a cena e conduz a sessão  & 1329 \\
\texttt{console.js}      & painel de informação dentro da cena           & 436 \\
\texttt{painel.js}       & painéis de texto desenhados em canvas         & 208 \\
\texttt{controles.js}    & controle de caminhada por toque               & 178 \\
\texttt{animador.js}     & reprodução das animações dos modelos          & 173 \\
\texttt{descoberta.js}   & portas de armário e revelação por aproximação & 136 \\
\midrule
\multicolumn{2}{l}{Total da experiência} & 3.378 \\
\bottomrule
\end{tabular}
\end{table}

O arquivo \texttt{dados.js} ocupa posição particular: ele é a fonte única
de verdade do sistema. Os cinco animais, os cinco níveis, as distâncias,
as quantidades e as escalas estão todos declarados ali, e nenhum outro
arquivo repete esses valores. Alterar a progressão dos níveis é alterar
uma linha --- o que importa num protótipo de pesquisa, em que os
parâmetros do estímulo podem precisar de ajuste entre sessões.

\subsection{Estímulos e ambientes}

Foram contemplados cinco animais frequentemente associados a fobias
específicas. Cada um é apresentado no ambiente em que seria encontrado,
construído com primitivas geométricas do próprio arcabouço.

\begin{table}[htbp]
\centering
\caption{Estímulos, fobias associadas e ambientes de apresentação.}
\label{tab:estimulos}
\begin{tabular}{llll}
\toprule
\textbf{Estímulo} & \textbf{Fobia} & \textbf{Ambiente} & \textbf{Tamanho real} \\
\midrule
Cobra  & Ofidiofobia     & Floresta & 1,40 m \\
Rato   & Musofobia       & Sala     & 0,25 m \\
Barata & Katsaridafobia  & Cozinha  & 0,05 m \\
Aranha & Aracnofobia     & Sala     & 0,08 m \\
Sapo   & Ranidafobia     & Quintal  & 0,10 m \\
\bottomrule
\end{tabular}
\end{table}

% [FIGURA] capturas dos quatro ambientes
A escolha de construir os cenários com primitivas, em vez de importar
ambientes prontos, decorre do peso: os quatro ambientes somados não
acrescentam nenhum arquivo ao projeto, o que mantém o tempo de
carregamento compatível com uma sessão de testes em que cada participante
abre a página do zero.

\subsubsection{Tamanho real e escala de apresentação}

Tamanho real e estímulo perceptível são grandezas distintas. Uma barata
de cinco centímetros, observada a três metros, ocupa menos de cinco
pixels na tela: fisicamente correta e visualmente inútil. A primeira
versão do protótipo apresentava os animais em tamanho real e o efeito
foi exatamente esse --- os estímulos pequenos eram imperceptíveis.

A solução adotada foi declarar separadamente, para cada animal, o tamanho
real e uma escala de apresentação. O sistema apresenta o produto dos
dois, e o trabalho pode informar ambos com transparência.

\begin{table}[htbp]
\centering
\caption{Tamanho real, escala de apresentação e tamanho apresentado.}
\label{tab:escalas}
\begin{tabular}{lrrr}
\toprule
\textbf{Estímulo} & \textbf{Real} & \textbf{Escala} & \textbf{Apresentado} \\
\midrule
Cobra  & 1,40 m & 1,0$\times$ & 1,40 m \\
Rato   & 0,25 m & 1,8$\times$ & 0,45 m \\
Barata & 0,05 m & 4,5$\times$ & 0,23 m \\
Aranha & 0,08 m & 3,6$\times$ & 0,29 m \\
Sapo   & 0,10 m & 3,0$\times$ & 0,30 m \\
\bottomrule
\end{tabular}
\end{table}

Note-se que a escala é inversamente proporcional ao tamanho do animal:
quanto menor o estímulo real, maior o fator necessário para torná-lo
perceptível. A cobra, única de porte comparável ao humano, é apresentada
em tamanho real.

\subsubsection{Ajuste automático dos modelos}

Modelos tridimensionais obtidos de repositórios públicos chegam em
unidades arbitrárias e com a origem em ponto imprevisível. Um mesmo
arquivo pode medir duas ou cento e vinte unidades. O sistema resolve isso
medindo a caixa envolvente do modelo depois do carregamento e calculando
o fator que o leva ao tamanho declarado, além de deslocá-lo
verticalmente para que a base encoste no piso.

Um terceiro ajuste foi necessário e é menos evidente: \emph{a maior
dimensão do modelo nem sempre é aquela a que o tamanho se refere}. O
modelo de cobra empregado inicialmente vinha em pose empinada, de modo
que a sua maior dimensão era a altura; escalar pela maior dimensão
produzia uma cobra de 1,2~m de altura, semelhante a um pilar. Cada animal
passou a declarar, portanto, a qual dimensão o seu tamanho se refere.

Um defeito dessa mesma natureza foi registrado durante o
desenvolvimento e serve de ilustração: a substituição do modelo de cobra
por outro, já em pose estendida, fez a configuração anterior escalar a
dimensão errada, e a verificação automatizada registrou estímulos de 13 a
16 metros de comprimento. O erro foi detectado pela ferramenta, e não a
olho nu.

\subsection{Níveis de exposição}

A progressão é organizada em cinco níveis, nos quais variam quatro
dimensões: a distância até o estímulo, a quantidade de exemplares, a
presença de movimento e a possibilidade de interação direta.

\begin{table}[htbp]
\centering
\caption{Dimensões que variam a cada nível de exposição.}
\label{tab:niveis}
\begin{tabular}{clccc}
\toprule
\textbf{Nível} & \textbf{Denominação} & \textbf{Distância} & \textbf{Movimento} & \textbf{Interação} \\
\midrule
1 & Primeiro contato      & máxima    & não & não \\
2 & Aproximação           & média     & não & não \\
3 & Movimento do estímulo & média     & sim & não \\
4 & Proximidade elevada   & reduzida  & sim & não \\
5 & Interação direta      & mínima    & sim & sim \\
\bottomrule
\end{tabular}
\end{table}

Os níveis 2 e 3 compartilham a mesma distância por definição: entre eles
muda o movimento, não a proximidade. Isolar uma variável por vez é o que
permite, na análise, atribuir eventual mudança de comportamento ao fator
correto.

\subsubsection{Distâncias}

Cada animal possui a sua própria escala de distâncias, porque os tamanhos
variam em uma ordem de grandeza. Cada um declara a distância do Nível~1 e
a do Nível~5; os intermediários são interpolados por uma proporção
declarada no nível.

\begin{table}[htbp]
\centering
\caption{Distância ao estímulo principal, em metros, por nível.}
\label{tab:distancias}
\begin{tabular}{lrrrrr}
\toprule
\textbf{Estímulo} & \textbf{N1} & \textbf{N2} & \textbf{N3} & \textbf{N4} & \textbf{N5} \\
\midrule
Cobra  & 5,50 & 3,61 & 3,61 & 2,14 & 1,30 \\
Rato   & 3,20 & 2,17 & 2,17 & 1,36 & 0,90 \\
Aranha & 2,60 & 1,81 & 1,81 & 1,20 & 0,85 \\
Sapo   & 2,60 & 1,81 & 1,81 & 1,20 & 0,85 \\
Barata & 1,90 & 1,38 & 1,38 & 0,98 & 0,75 \\
\bottomrule
\end{tabular}
\end{table}

\subsubsection{Quantidade}

A quantidade de exemplares cresce de um até quinze, ajustada por um fator
próprio de cada espécie. A necessidade do fator veio da observação: uma
cozinha com vinte e quatro baratas é uma cena reconhecível; uma mata com
quinze cobras de 1,40~m não é --- lê-se como exagero, e não como
intensificação do estímulo.

\begin{table}[htbp]
\centering
\caption{Número de exemplares apresentados por nível.}
\label{tab:quantidades}
\begin{tabular}{lrrrrrr}
\toprule
\textbf{Estímulo} & \textbf{Fator} & \textbf{N1} & \textbf{N2} & \textbf{N3} & \textbf{N4} & \textbf{N5} \\
\midrule
Barata & 1,6 & 1 & 5 & 10 & 16 & 24 \\
Rato   & 1,0 & 1 & 3 & 6  & 10 & 15 \\
Aranha & 1,0 & 1 & 3 & 6  & 10 & 15 \\
Sapo   & 1,0 & 1 & 3 & 6  & 10 & 15 \\
Cobra  & 0,5 & 1 & 2 & 3  & 5  & 8  \\
\bottomrule
\end{tabular}
\end{table}

O Nível~1 ignora o fator: o primeiro contato é sempre com um único
exemplar, para todos os participantes, de modo que seja idêntico e
previsível.

\subsubsection{Reprodutibilidade da disposição}

As posições dos exemplares são sorteadas, mas o sorteio parte de uma
semente derivada do animal e do nível. A consequência é que o mesmo nível
monta sempre a mesma cena. Sem isso, dois participantes veriam
disposições diferentes e a comparação entre eles perderia sentido --- um
poderia ter encontrado três animais a um metro e o outro, nenhum.

\subsubsection{Exclusão deliberada das animações de ataque}

Os modelos empregados trazem animações de ataque e de morte, nomeadas
\texttt{Attack} e \texttt{Death}. Nenhuma delas é utilizada. Um animal
que ataca produz susto, e susto é o oposto de exposição gradual e
controlada: retira do participante o domínio sobre o próprio ritmo, que é
o princípio do procedimento.

A exclusão é uma decisão de projeto fundamentada, e não uma limitação
técnica --- as animações estão nos arquivos e poderiam ser reproduzidas
com uma linha de configuração.

\subsection{Comportamento dos estímulos}

Um grupo em que todos os exemplares se movem de maneira idêntica não
parece um grupo: parece um enfeite. Cada exemplar recebe, portanto, um
comportamento próprio, sorteado com a mesma semente que determina as
posições.

\begin{table}[htbp]
\centering
\caption{Comportamentos atribuídos a cada exemplar.}
\label{tab:comportamentos}
\small
\begin{tabular}{p{2.9cm}p{6.4cm}p{2.8cm}}
\toprule
\textbf{Comportamento} & \textbf{O que faz} & \textbf{Custo} \\
\midrule
Quieto    & apenas respira, com balanço vertical lento & baixo: dispensa o reprodutor de animação \\
Inquieto  & reproduz a animação do arquivo e gira o corpo & médio \\
Andarilho & o mesmo, e ainda caminha devagar entre dois pontos & maior \\
\bottomrule
\end{tabular}
\end{table}

Nos níveis 1 e 2, 65\% dos exemplares permanecem quietos; a partir do
Nível~3, a proporção passa a 25\% quietos, 40\% inquietos e 35\%
andarilhos. Três regras são mantidas em todos os casos:

\begin{enumerate}
  \item \textbf{nenhum exemplar fica completamente imóvel} --- quem não
  reproduz animação recebe ao menos o balanço de respiração;
  \item \textbf{sempre há algum exemplar parado} --- sem contraste, o
  grupo inteiro se torna um borrão de movimento e nada chama atenção;
  \item \textbf{o exemplar principal nunca se desloca} --- ele está na
  distância exata que define o nível, e sair do lugar desfaria a
  progressão que se pretende medir.
\end{enumerate}

\subsection{Exploração e descoberta}

A partir do Nível~2, parte dos exemplares inicia a cena invisível,
alojada em esconderijos distribuídos pelos ambientes. São vinte e um
esconderijos no total --- cinco na sala, nove na cozinha, cinco na
floresta e cinco no quintal --- e a fração de exemplares escondidos é de
35\% nos níveis 2 e 3 e de 45\% nos níveis 4 e 5.

Dois modos de revelação foram implementados:

\begin{table}[htbp]
\centering
\caption{Modos de revelação dos exemplares escondidos.}
\label{tab:revelacao}
\begin{tabular}{p{2.6cm}p{10.0cm}}
\toprule
\textbf{Esconderijo} & \textbf{Como o exemplar aparece} \\
\midrule
Com porta & somente quando o participante \textbf{abre a porta do armário}; a aproximação não basta \\
Sem porta & quando o participante \textbf{chega a menos de 1,6~m} do esconderijo \\
\bottomrule
\end{tabular}
\end{table}

% [FIGURA] cozinha com as portas fechadas e abertas
Dos nove esconderijos da cozinha, cinco são armários com porta articulada
que gira ao ser tocada. Ao aparecer, o exemplar desloca-se para fora do
esconderijo, em direção ao centro do ambiente.

O efeito pretendido é deslocar o controle da exposição. Em lugar de
encarar todos os estímulos de uma só vez, o participante avança no
próprio ritmo e decide quando abrir a próxima porta. O painel de
informação indica quantos exemplares já foram encontrados do total
escondido --- sem esse indicador, quem explora tende a desistir supondo
já ter visto tudo.

\subsection{Interface}
\label{sec:interface}

A interface da aplicação é construída como objeto da cena
tridimensional, e não como camada de HTML sobreposta ao desenho. A
decisão tem motivação técnica direta: ao entrar em modo de Realidade
Virtual, o navegador exibe exclusivamente o conteúdo tridimensional, de
modo que qualquer elemento em HTML desaparece. Um painel sobreposto
deixaria o participante de óculos sem nível, sem instrução e sem
controles.

% [FIGURA] o console dentro da cena
O painel acompanha o corpo do participante, e não a sua cabeça:
desloca-se com ele quando caminha, mas não gira junto com o olhar.
Interface fixada ao campo de visão é reconhecidamente causa de
desconforto em Realidade Virtual.

\subsubsection{Dois arranjos, determinados pelo campo de visão}

O posicionamento da interface precisou ser calculado, e não estimado. A
câmera do arcabouço mede os seus 80 graus de abertura na
\emph{vertical}; a abertura horizontal resulta da proporção da tela.
Daí decorre uma diferença considerável:

\begin{table}[htbp]
\centering
\caption{Campo de visão horizontal conforme o formato da tela.}
\label{tab:fov}
\begin{tabular}{lrr}
\toprule
\textbf{Tela} & \textbf{Proporção} & \textbf{Abertura horizontal} \\
\midrule
Computador (1280 $\times$ 800) & 1,60 & 106$^\circ$ \\
Celular em pé (390 $\times$ 844) & 0,46 & 42$^\circ$ \\
\bottomrule
\end{tabular}
\end{table}

Com 42 graus de abertura, metade é 21 graus --- e um painel posicionado a
32 graus à esquerda, adequado no computador, fica \emph{inteiramente fora
da tela} em um celular. A aplicação adota, por isso, dois arranjos: no
computador e no óculos o painel fica à esquerda, com os botões abaixo
dele; no celular tudo se centraliza e encolhe, aproveitando a altura que
a tela em pé tem de sobra.

\subsubsection{Dois apontadores}

A seleção de alvos revelou um conflito que merece registro, por ser
próprio de interfaces tridimensionais e não ter equivalente em páginas
convencionais.

Em celulares sem controle de mão, a seleção se dá pelo olhar: a mira
permanece no centro da tela e, após permanecer sobre um alvo por 1,2
segundo, aciona-o. Como no celular o painel de interface é centralizado,
a mira repousava sobre ele e acionava sozinha o botão que estivesse no
centro --- se fosse o de reiniciar, a página recarregava, e recarregava
novamente, em laço.

A correção foi separar os apontadores conforme a sua natureza: a mira
enxerga apenas os estímulos e as portas, enquanto um segundo apontador,
que parte do dedo ou do cursor e nunca do centro da tela, enxerga apenas
a interface. Dentro do óculos, onde não há dedo nem cursor, a mira volta
a enxergar ambos.

Um segundo defeito foi detectado na mesma ocasião: o arcabouço não ignora
objetos invisíveis ao calcular o alcance do raio de seleção, de modo que
um painel oculto permanecia interceptando cliques destinados aos botões
atrás dele. A solução foi distinguir o que \emph{é} interface do que
\emph{pode ser acionado no momento}, atualizando a lista de alvos a cada
abertura ou fechamento de painel.

Ambos os defeitos só se manifestaram quando os testes automatizados
passaram a empregar cliques e toques reais. Até então disparavam o evento
diretamente sobre o elemento, o que contorna o cálculo do raio --- e
teria aprovado os dois casos sem apontar problema. O registro fica como
observação de método: em interface tridimensional, teste por evento
sintético produz falso positivo.

\subsection{Deslocamento}

\begin{table}[htbp]
\centering
\caption{Formas de interação conforme o dispositivo.}
\label{tab:controles}
\small
\begin{tabular}{p{2.2cm}p{2.6cm}p{3.5cm}p{3.6cm}}
\toprule
\textbf{Dispositivo} & \textbf{Olhar} & \textbf{Caminhar} & \textbf{Selecionar} \\
\midrule
Computador & arrastar o cursor & teclas W A S D & clique \\
Celular & girar o aparelho & controle circular na tela & toque \\
Óculos & mover a cabeça & deslocamento físico & raio do controle ou olhar \\
\bottomrule
\end{tabular}
\end{table}

No celular não há teclado, e girar o aparelho permite olhar em volta mas
não mover-se --- o que inviabilizaria tanto a exploração quanto a
descoberta dos exemplares escondidos. Foi implementado, para esse caso,
um controle circular de toque, cujo deslocamento é \emph{relativo ao
olhar} e limitado pelas paredes do cenário, com folga de 45~cm.

A velocidade adotada é de 1,4~m/s, próxima da caminhada humana. Valores
superiores produzem desconforto em Realidade Virtual.

\subsection{Verificação técnica}
\label{sec:verificacao}

% [FIGURA] a ferramenta de verificação em execução
O protótipo conta com ferramenta própria de verificação automatizada,
que percorre as vinte e cinco combinações de animal e nível e confere, em
cada uma: o carregamento do modelo, a escala aplicada, o apoio no piso, o
enquadramento no campo de visão, o carregamento do áudio, a taxa de
quadros sustentada e a ausência de erros de execução.

A ferramenta não é um acessório. Foi ela que detectou o erro de escala da
cobra mencionado anteriormente, e é ela que permite afirmar, com número,
que uma alteração não degradou o que já funcionava.

\subsubsection{Desempenho}

A medição é feita em renderização por software, sem aceleração por placa
de vídeo --- condição deliberadamente desfavorável. Em qualquer aparelho
com placa de vídeo o desempenho tende a atingir o limite da tela.

A Tabela~\ref{tab:desempenho} compara o desempenho antes e depois do
conjunto de otimizações descrito a seguir, na mesma máquina e com a mesma
área renderizada.

\begin{table}[htbp]
\centering
\caption{Quadros por segundo antes e depois das otimizações, com o
número de exemplares entre parênteses.}
\label{tab:desempenho}
\begin{tabular}{llrrrrr}
\toprule
\textbf{Estímulo} & & \textbf{N1} & \textbf{N2} & \textbf{N3} & \textbf{N4} & \textbf{N5} \\
\midrule
Cobra  & antes  & 16 (1) & 11 (3) & 8 (6)  & 6 (10)  & 4 (15) \\
       & depois & 21 (1) & 20 (2) & 19 (3) & 18 (5)  & 16 (8) \\
\addlinespace
Rato   & antes  & 43 (1) & 35 (3) & 26 (6) & 20 (10) & 15 (15) \\
       & depois & 46 (1) & 40 (3) & 31 (6) & 30 (10) & 21 (15) \\
\addlinespace
Barata & antes  & 37 (1) & 35 (3) & 34 (6)  & 32 (10) & 29 (15) \\
       & depois & 37 (1) & 36 (5) & 34 (10) & 32 (16) & 29 (24) \\
\addlinespace
Aranha & antes  & 44 (1) & 37 (3) & 29 (6) & 24 (10) & 19 (15) \\
       & depois & 47 (1) & 41 (3) & 35 (6) & 32 (10) & 25 (15) \\
\addlinespace
Sapo   & antes  & 26 (1) & 22 (3) & 18 (6) & 14 (10) & 11 (15) \\
       & depois & 26 (1) & 24 (3) & 21 (6) & 19 (10) & 14 (15) \\
\bottomrule
\end{tabular}
\end{table}

A média das vinte e cinco combinações passou de 24 para 29 quadros por
segundo, e o número de combinações sem aviso de desempenho subiu de oito
para doze. Em nenhuma das duas medições houve falha.

Três leituras merecem registro:

\begin{itemize}
  \item \textbf{a cozinha passou de 15 para 24 exemplares sem perder
  desempenho} (29 quadros por segundo em ambas as medições, no Nível~5).
  O que paga esse aumento é o comportamento \emph{quieto}, que dispensa
  o reprodutor de animação, somado ao limite de oito exemplares animados
  simultaneamente;

  \item \textbf{a floresta passou de 4 para 16 quadros por segundo no
  Nível~5}, por duas causas somadas: a redução das texturas do modelo de
  cobra de 1024 para 512 pixels e a redução pela metade da quantidade de
  exemplares;

  \item \textbf{a mesma quantidade tornou-se mais barata}: com quinze
  ratos, de 15 para 21 quadros por segundo; com quinze aranhas, de 19
  para 25. Nada mudou no modelo nem no cenário --- a diferença está
  inteiramente no comportamento dos exemplares.
\end{itemize}

\subsubsection{Otimizações aplicadas}

\begin{table}[htbp]
\centering
\caption{Otimizações e respectivo efeito.}
\label{tab:otimizacoes}
\begin{tabular}{p{4.2cm}p{8.4cm}}
\toprule
\textbf{Otimização} & \textbf{Efeito} \\
\midrule
Piso pintado em canvas &
o padrão do piso era composto por 90 peças separadas; passou a ser
pintado uma vez numa textura e repetido: um desenho em lugar de noventa \\
Céu como cor de fundo &
substituição da esfera de céu por cor de fundo da cena \\
Limite de exemplares animados &
no máximo oito reproduzem a animação do arquivo; cada um recalcula a
posição de cada osso a cada quadro \\
Comportamento \emph{quieto} &
exemplares parados não recebem reprodutor de animação \\
Redução de texturas &
cobra: cinco texturas de 1024 para 512 pixels, de 6,13~MB para 3,82~MB \\
Painel em canvas único &
a interface inteira é desenhada numa só textura: um objeto por quadro,
em lugar de nove \\
\bottomrule
\end{tabular}
\end{table}

\subsection{Modelos tridimensionais}

\begin{table}[htbp]
\centering
\caption{Modelos empregados e animações utilizadas.}
\label{tab:modelos}
\begin{tabular}{lrrll}
\toprule
\textbf{Arquivo} & \textbf{Tamanho} & \textbf{Triângulos} & \textbf{Usadas} & \textbf{Descartadas} \\
\midrule
cobra.glb  & 3,82 MB & 4.380 & Animation   & --- \\
aranha.glb & 0,43 MB & 2.712 & Idle, Walk, Jump & Attack, Death \\
sapo.glb   & 0,56 MB & 4.920 & Idle, Jump  & Attack, Death \\
rato.glb   & 0,53 MB & 4.004 & Idle, Walk, Jump & Run, Attack, Death \\
barata.glb & 0,52 MB & 576   & idle, walking & Attack, death, flying \\
\bottomrule
\end{tabular}
\end{table}

A reprodução das animações exigiu um componente próprio: o arcabouço
carrega o modelo e as suas animações, mas não as reproduz. O componente
desenvolvido faz a ponte com o reprodutor do Three.js e localiza o clipe
por correspondência parcial de nome --- um modelo pode nomear a sua
animação de repouso como \texttt{SnakeArmature|Snake\_Idle}, e procurar
por \texttt{Idle} já a encontra. Isso evita depender do nome exato, que
varia de modelo para modelo.

\subsection{Áudio}

Cada animal possui um som próprio, posicionado na cena, de modo que o
volume percebido varia com a distância --- a progressão dos níveis também
se ouve. Há ainda um som de ambiente, não posicionado.

O áudio começa desligado e depende de ativação explícita. Não se trata de
opção de interface: navegadores proíbem a reprodução de áudio antes de
uma interação do usuário com a página, e como o participante chega vindo
de outra página, a interação anterior não conta. A preferência é
preservada entre níveis durante a mesma sessão.

\subsection{Limites conhecidos do artefato}

Registram-se, por honestidade metodológica, as limitações do artefato que
são de natureza técnica:

\begin{itemize}
  \item os estímulos pequenos são apresentados em escala maior que a
  real, pelas razões expostas; o fator está documentado, mas a
  apresentação não é fisicamente fiel;

  \item os sons foram sintetizados e não são gravações dos animais;

  \item não há medida fisiológica. A possibilidade de captar frequência
  cardíaca foi investigada e descartada: exige sensor com perfil padrão
  de transmissão, que os relógios de uso comum não expõem, e
  acrescentaria ao estudo uma natureza clínica que o seu escopo não
  comporta;

  \item o desempenho medido refere-se à renderização por software;
  números obtidos no aparelho efetivamente utilizado nas sessões devem
  ser coletados com a mesma ferramenta.
\end{itemize}
